% Source: https://artofproblemsolving.com/community/c3h1695826_parallel
% Stable identity: aops-topic-1695826-q01
The cevians $AM$, $BF$, $CE$ all pass through $P$, so by Ceva's theorem
\[
\dfrac{AE}{EB}\cdot\dfrac{BM}{MC}\cdot\dfrac{CF}{FA}=1 .
\]
Since $M$ is the midpoint of $BC$, $\dfrac{BM}{MC}=1$, hence $\dfrac{AE}{EB}=\dfrac{FA}{CF}$, i.e.\ $\dfrac{AE}{EB}=\dfrac{AF}{FC}$. Adding $1$ to both sides gives $\dfrac{AB}{EB}=\dfrac{AC}{FC}$, so $\dfrac{AE}{AB}=\dfrac{AF}{AC}$. Then triangles $AEF$ and $ABC$ share the angle at $A$ and have proportional sides about it, so they are similar (SAS), $\angle AEF=\angle ABC$, and by corresponding angles $\boxed{EF\parallel BC}$. $\blacksquare$

\begin{asy}
usepackage("amsmath");
size(6cm);
defaultpen(fontsize(10pt));
pair A=(1.3,4),B=(0,0),C=(6,0),M=(B+C)/2;
pair P=A+0.55*(M-A);
pair Ep=extension(C,P,A,B),F=extension(B,P,A,C);
draw(A--B--C--cycle,linewidth(0.9));
draw(A--M,linewidth(0.7));
draw(C--Ep,linewidth(0.6));draw(B--F,linewidth(0.6));
draw(Ep--F,red+dashed+linewidth(0.8));
dot(P);dot(M);dot(Ep);dot(F);
label("$A$",A,N);label("$B$",B,SW);label("$C$",C,SE);label("$M$",M,S);
label("$P$",P,dir(200));label("$E$",Ep,W);label("$F$",F,NE);
label("\mbox{}",(min(currentpicture,user=true)+max(currentpicture,user=true))/2,black); // keeps the text after the figure black when LaTeX typesets the labels inline
\end{asy}
